% CS 423/520 Computer Architecture, Fall 2026: Bonus HW, Critical Paper Review
%
% How to use this template
%   1. Fill in the five fields below (name, student ID, paper key, title, venue).
%   2. Replace every placeholder text in the sections with your review.
%   3. Compile with pdflatex (or upload this file to Overleaf) and check that
%      the PDF has at most two pages.
%   4. Name the PDF  review_<studentID>_<paperkey>.pdf  and upload it to Moodle.
%
% Do NOT change anything above the line "FILL IN BELOW": no other font, font
% size, margins, line spacing, or page size. Do NOT remove or rename sections.
% A review longer than two pages gets 0 points.
\documentclass[11pt,a4paper]{article}

\usepackage[margin=2.5cm]{geometry}
\usepackage[T1]{fontenc}
\usepackage{lmodern}
\usepackage{enumitem}
\usepackage{xcolor}
\usepackage{fancyhdr}
\usepackage{lastpage}
\usepackage[hidelinks]{hyperref}

\setlength{\parindent}{0pt}
\setlength{\parskip}{3pt}
\setlist{nosep,leftmargin=2.2em}

% Section headings: fixed titles, compact spacing.
\makeatletter
\renewcommand\section{\@startsection{section}{1}{0pt}%
  {6pt plus 2pt}{2pt}{\large\bfseries}}
\renewcommand\subsection{\@startsection{subsection}{2}{0pt}%
  {3pt plus 1pt}{1pt}{\normalsize\bfseries}}
\makeatother
\newcommand{\summaryitem}[1]{\par\smallskip\textbf{#1.}\ }

\newcommand{\reviewer}[1]{\def\theReviewer{#1}}
\newcommand{\studentid}[1]{\def\theStudentID{#1}}
\newcommand{\paperkey}[1]{\def\thePaperKey{#1}}
\newcommand{\papertitle}[1]{\def\thePaperTitle{#1}}
\newcommand{\papervenue}[1]{\def\thePaperVenue{#1}}

\pagestyle{fancy}
\fancyhf{}
\renewcommand{\headrulewidth}{0.4pt}
\fancyhead[L]{\small CS 423/520 Bonus HW: Paper Review}
\fancyhead[R]{\small\texttt{\thePaperKey}\quad\theStudentID}
\fancyfoot[C]{\small\thepage/\pageref*{LastPage}}

% Visible warning when the review is longer than two pages.
\AtEndDocument{%
  \ifnum\value{page}>2
    \par\bigskip\noindent\fcolorbox{red}{red!8}{\parbox{\dimexpr\linewidth-2\fboxsep-2\fboxrule}{%
      \color{red}\bfseries This review is longer than two pages. Reviews longer
      than two pages get 0 points. Shorten it before you submit.}}%
  \fi}

% ============================== FILL IN BELOW ===============================
\reviewer{Your Name}
\studentid{XXXXXXXX}
\paperkey{e.g., kim2014flipping}   % the key of the paper, exactly as on the paper list
\papertitle{The full title of the paper that you review}
\papervenue{ISCA 2014}       % venue and year, as on the paper list
% ============================================================================

\begin{document}

\begin{center}
  {\Large\bfseries \thePaperTitle}\\[2pt]
  {\small \thePaperVenue\quad$\bullet$\quad Paper key: \texttt{\thePaperKey}}\\[2pt]
  {\small Reviewed by \theReviewer\ (\theStudentID)}
\end{center}

% ---------------------------------------------------------------------------
\section*{Summary}
% Write the summary in your own words, in this order. Do not copy sentences
% from the paper. One to three sentences per item is usually enough.

\summaryitem{Problem \& Goal}
What problem does the paper solve? Why is it important? What does the paper want to achieve?

\summaryitem{Key Ideas/Solution}
What is the main idea of the solution, in one or two sentences?

\summaryitem{Novelty}
What is new in this paper? What did nobody do before?

\summaryitem{Mechanisms \& Implementation}
How does the solution work? What hardware or software does it need?

\summaryitem{Major Results}
What are the most important results? Give numbers and say what they are compared to.

\summaryitem{Takeaways/Conclusions}
What should a reader remember from this paper?

\summaryitem{State of the Art Before This Paper}
What were the best solutions before this paper, and why were they not good enough?

% ---------------------------------------------------------------------------
\section*{Strengths}
% The most important strengths, in a short list. One or two sentences each.
% At least three.
\begin{enumerate}
  \item First strength.
  \item Second strength.
  \item Third strength.
\end{enumerate}

% ---------------------------------------------------------------------------
\section*{Weaknesses}
% The most important weaknesses, in a short list. One or two sentences each.
% At least three.
\begin{enumerate}
  \item First weakness.
  \item Second weakness.
  \item Third weakness.
\end{enumerate}

% ---------------------------------------------------------------------------
\section*{Comments for the Authors}
% One paragraph for each weakness above, with the same number (Weakness 1, 2, ...).
% Define the weakness clearly and precisely, and give an example.
\subsection*{Weakness 1}
Explain the first weakness: what exactly is missing or wrong, why it matters, and an example.
\subsection*{Weakness 2}
Explain the second weakness: what exactly is missing or wrong, why it matters, and an example.
\subsection*{Weakness 3}
Explain the third weakness: what exactly is missing or wrong, why it matters, and an example.


% ---------------------------------------------------------------------------
\section*{Thoughts/Ideas}
% Your own thinking. Answer the questions that matter most for this paper,
% for example: What did you learn, enjoy, or dislike, and why? Can you do
% better? How would you eliminate the weaknesses you found? Do you have a
% fundamentally different solution to the same problem? Does the paper bring
% new problems, and how would you solve them?
Your thoughts and ideas. Include at least one concrete idea: what you would change, and how you would test whether it works.

% ---------------------------------------------------------------------------
\section*{GenAI Statement}
% One line. Examples: "I did not use GenAI tools." or
% "I used <tool> only to proofread the final text."
I did not use GenAI tools.

\end{document}
